\section{Background: LFM2.5-VL}
\label{sec:background}

LFM2.5-VL-1.6B~\cite{lfm2} consists of three components.

\paragraph{Vision encoder.} A SigLIP2 NaFlex so400m ViT~\cite{siglip2} (27 layers, hidden size 1152,
16 heads, patch size 16, $\sim$400M parameters). Images up to $512\times512$ are processed at native
aspect ratio. Larger images are split into 2--10 tiles of $512\times512$ plus a thumbnail.

\paragraph{Projector.} A $2\times2$ pixel-unshuffle (reducing image tokens by $4\times$) followed by a
two-layer MLP ($4608 \rightarrow 2048 \rightarrow 2048$, GELU), which gives 64--256 tokens per tile.

\paragraph{Language model.} LFM2.5-1.2B: 16 layers, of which 10 are \emph{gated short-convolution}
blocks and 6 are grouped-query attention blocks (32 query / 8 KV heads, head size 64, QK-RMSNorm, RoPE).
Every layer is followed by a SwiGLU MLP (feed-forward width 8192), with pre-RMSNorm and residual
connections throughout. The gated short-convolution block computes
\begin{equation}
  [B, C, x] = W_\text{in} h,\qquad
  u = \operatorname{conv1d}_{k=3}(B \odot x),\qquad
  y = W_\text{out}(C \odot u),
\end{equation}
where $\operatorname{conv1d}_{k=3}$ is a causal depthwise convolution. During decoding it only needs a
constant-size state per layer. Because only 6 layers use attention, the KV cache is $\sim$12\,KB per token.

\todo{Figure: architecture diagram (figures/architecture.pdf).}

\paragraph{Target hardware.} NVIDIA RTX PRO 5000 Blackwell (compute capability 12.0).
\todo{Confirm SM count, memory size/bandwidth, TDP, clocks from nvidia-smi / deviceQuery.}
